\begin{frame}[t]{\ev{\tagM\,\tagP}The observed gaps define obligations that future branches must preserve.}
\begin{center}\lifecycle{0}\end{center}
\vspace{.1cm}
\begin{ticks}
\item \textbf{Observed needs:} explicit tests, accurate receipts, durable export.
\item \textbf{Proposed branches:} identified parent, private continuations, fresh authority.
\item This case establishes files and feedback, without a demonstrated need to preserve live processes.
\end{ticks}
\sources{SELFHOST-2/3 observations; the integrated branching lifecycle is proposed.}
\qadetail{Before the first repair, we could identify the input tree, dependencies, and failed-test receipt for proposed alternative continuations. No actual fork occurred. The case does not establish that useful process memory must survive, so worktrees or warm environments remain alternatives. The hypothetical live-application failure is a future workload, not factory evidence. New scope needs authorization and a behavioral obligation; missing coverage needs a targeted artifact-bound test; inaccurate receipts need typed event semantics; recovery needs export before teardown; replay needs retrievable inputs. For proposed branching, declare the captured surface, sharing policy, external observations, child identity, and authority. The numerical reducer does not detect the missing test by itself.}
\note{[0:55] We can now connect each observation to a required mechanism. A new behavior needs a concrete test obligation. Records need the actual execution status and actor. Candidates need durable export on every exit. For future branches, we also need an identified parent, private changes, and fresh authority. We could freeze the pre-repair files and feedback for alternatives, but this case does not show that live process memory must survive. The branching lifecycle remains proposed. We next compare the mechanisms that preserve each kind of required state.}
\end{frame}
